% 本文件由 LaTeX 排版 Agent 根据有效 translation.json 自动生成。
% author=Jerome; work_id=3006; title=被掳隐修士玛耳苛传

\chapter{被掳隐修士\Person{玛耳苛}传}

% structure: p0001 source=h1 target=omitted reason=page-title-already-rendered-by-parent

\Emph{玛耳苛传写于\Place{白冷}，公元391年。其来源与目的在第一、二章已充分说明。}\par

凡准备海战的人，先在海港和静水中练习掌握船舵、划动船桨、准备铁钩和搭钩。他们把士兵列在甲板上，使他们习惯于在滑动的地面上站稳脚步；这样，当真战来临时，他们就不会因曾在模拟战中经历过而退缩。我的情形也是如此。我长久保持沉默，因为一位我提及他时便使他痛苦的人对我施压。但现在，我准备撰写更广泛的历史，我愿意通过这部小作品练习自己，仿佛擦去舌上的锈迹。因为我立意（若天主赐我寿命，并且责备我的人最终停止迫害我这逃亡者并幽居隐修院的人）写一部基督教会史，从我们救主的来临直到我们自己的时代，即从宗徒直到我们生活的末叶，并显明它如何诞生、藉什么工具，以及它如何随迫害而成长、以殉道而加冕；然后，在信奉基督的皇帝之后，它如何在影响和财富上增长，而在基督徒美德上却减退。但这些留待别处。现在回到正题。\par

\Place{玛洛尼亚}是\Place{叙利亚}\Place{安提约基雅}以东约三十英里的一座小村庄。它曾有过许多主人或地主，当我年轻时客居叙利亚时，它落入我的密友\Person{厄瓦格略}主教之手，我现在提及他的名字是为了说明我的信息来源。那地方当时有一位名叫\Person{玛耳苛}的老人，他的名字可以译为\Quote{王}，他是叙利亚人，语言也是叙利亚语，事实上是地地道道的本地人。他的同伴是一位老妇人，非常衰老，似乎行将就木；两人都如此热心虔诚，并经常出入教会，人们会以为他们是\BibleBook{}{福音}中的\Person{匝加利亚}和\Person{依撒伯尔}，只是不见\Person{若翰}。我怀着好奇问邻居们，他们之间是什么关系；是婚姻，是亲属，还是圣神的纽带？众人异口同声回答说他们是圣人，中悦天主，并给我讲了一个关于他们的奇异故事。我渴望知道更多，便急切地开始询问那人我所听到的是否属实，并得知如下。\par

我儿，他说，我曾在\Place{尼西比斯}耕种一小块地，且是独子。我的父母视我为继承人和他们家族唯一的幸存者，想强迫我结婚，但我拒绝说，我宁愿做隐修士。我父亲如何威胁、母亲如何哄骗我出卖我的贞洁，无需其他证据，只凭我逃离家庭和父母这一事实。我不能往东去，因为波斯就在附近，边境有罗马士兵把守；因此我转向西行，随身带了些旅途口粮，但仅够免于匮乏。简而言之，我最终来到\Place{卡尔基斯}旷野，它位于\Place{伊玛}和更南方的\Place{贝若亚}之间。在那里我找到一些隐修士，便将自己置于他们的指导下，靠双手劳动挣取生计，并以斋戒克制肉体的放纵。多年后，我心里起了回乡的愿望，想与母亲同住，在她有生之年安慰她的寡居（因为父亲，据我所知，已经去世），然后卖掉那小份产业，一部分施舍给穷人，一部分捐给隐修院（我羞愧地承认自己的不忠），自己留一些用于安逸的生活。我的院长开始大声说，这是魔鬼的诱惑，在貌似正当的借口下隐藏着那古老仇敌的陷阱。他说，这就像狗转回吃它所吐的。许多隐修士曾被这样的建议欺骗，因为魔鬼从不公开显现。他向我列举了\WorkTitle{圣经}中的许多例子，并告诉我甚至起初的\Person{亚当}和\Person{厄娃}也是被他用成为神祇的盼望所推翻的。当他未能说服我时，便跪倒在地，恳求我不要离弃他，也不要在手扶着犁后向后看而毁灭自己。不幸的是，我竟不幸战胜了我的劝告者。我以为他所寻求的不是我的得救，而是他自己的安逸。于是他送我离开隐修院，仿佛是去参加葬礼，最后向我告别说：\Quote{我看你带着\Person{撒旦}之子的印记。我不问你理由，也不接受你的借口。离群之羊立刻成为狼的口中之食。}\par

从\Place{贝若亚}往\Place{厄德撒}的路上，在公路旁是一片荒地，撒拉逊人在那里来回游荡，没有固定住处。旅客们因惧怕他们而成群结队同行，以便互相帮助摆脱迫在眉睫的危险。与我同行的有男人、女人、老人、青年、孩子，总共约七十人。突然间，依市玛耳人骑着马和骆驼袭击了我们，他们头发蓬松，束着头带，身体半裸，穿着宽大的军靴，斗篷在身后飘动，肩上挂着箭袋。他们拿着未上弦的弓，挥着长矛；因为他们来不是为了战斗，而是为了掠夺。我们被捉住，被分散，被带往不同方向。与此同时，我后悔自己的决定已晚，远未得到我的产业，而是与另一名可怜的受害者——一个女人——被分配给同一位主人使唤。我们被领着，更不如说是被高高地驼在骆驼背上，穿越一片荒芜的旷野，时刻预期死亡，与其说是坐着，不如说是悬着。半生不熟的肉是我们的食物，骆驼奶是我们的饮料。\par

5. 最后，我们渡过一条大河，来到沙漠腹地。按照当地习俗，我被命令向女主人和她的孩子们致敬，我们低下了头。在这里，我像一个囚徒一样换了衣服，也就是学会赤裸行走，因为天气酷热，除了腰间的遮盖物外不允许穿任何衣服。有人给了我一些羊去放牧，相对来说，我发现这个职业是一种安慰，因为我很少见到主人或同伴奴隶。我的命运似乎像圣经历史中的雅各，也让我想起梅瑟；两人都曾在沙漠中当过牧羊人。我以新鲜的奶酪和牛奶为食，不断祈祷，并唱我在隐修院学到的圣咏。我对被掳感到高兴，并感谢天主，因为我在沙漠中找到了我差点在自己的国家失去的隐修生活。\par

6. 但是，任何境况都无法阻挡魔鬼。他的陷阱何其多，难以言表！虽然我隐藏着，但他的恶意还是找到了我。我的主人看到羊群增长，发现我没有欺诈（我知道宗徒曾命令，主人应如同事奉天主一样忠实地服事），为了奖赏我以确保我更忠诚，就把那个曾与我一同被掳的女人给了我。我拒绝并说我是基督徒，而且只要她的丈夫还活着（她的丈夫与我们一同被掳，但被另一个主人带走了），我娶她是不合法的。我的主人毫不留情，怒火中烧，拔出剑来开始向我砍来。如果我没有立刻伸出手接手那个女人，他当场就会杀了我。唉！这时，一个比往常更黑的夜晚来临了，对我来说，来得太快了。我把新娘领进一个旧山洞；悲伤是伴娘；我们彼此退缩，但没有承认。那时我才真正感受到被掳；我趴在地上，开始哀叹我失去的隐修生活，并说：\Quote{我真不幸！我竟被保存到这种地步吗？我的邪恶把我带到了这个地步，让我在白发之年失去童贞，成为一个已婚男人？我若做了我本想避免的事，那为主而轻视父母、家乡、财产又有什么用呢？然而，也许我落到这个地步，是因为我思念家乡。我们该怎么办，我的灵魂？我们是毁灭，还是得胜？我们是等候主的手，还是用我们自己的剑自杀？把剑对准你自己；我的灵魂，我必须害怕你的死亡，胜于身体的死亡。保持的贞洁有自己的殉道。让为基督作证的人无人埋葬在沙漠；我将同时是迫害者和殉道者。}说着，我拔出剑来，即使在黑暗中也闪闪发光，把剑尖对准自己，说：\Quote{别了，不幸的女人：把我当作殉道者，而不是丈夫。}她扑倒在我脚前，喊道：\Quote{我因耶稣基督恳求你，并因这考验的时刻命令你，不要流你的血，把罪咎加在我身上。如果你选择死，先把你的剑对准我。让我们在这些条件下结合。假如我的丈夫回到我身边，我也将保持我在被掳中学到的贞洁；我宁愿死也不愿失去它。你为什么为了阻止与我的结合而死？如果你愿意，我就去死。把我当作你贞洁的伴侣吧；在精神的结合中爱我，胜过仅仅是肉体的爱。让我们的主人相信你是我的丈夫。基督知道你是我的兄弟。我们很容易使他们相信我们结婚了，因为他们看到我们如此相爱。}我承认，我惊异不已，之前就很钦佩这个女人的美德，现在我更把她当作妻子来爱。然而，我从未凝视过她的裸体；我从未碰过她的肉体，因为我害怕在和平中失去我在冲突中保全的东西。在这种奇怪的婚姻中，许多天过去了。婚姻使我们在主人面前更讨人喜欢，而且他们对我们逃跑没有怀疑；有时我甚至整个月都不在，像一个忠诚的牧羊人穿越荒野。\par

7. 过了很长一段时间，有一天我独自坐在沙漠中，除了大地和天空什么也看不见，我开始飞快地思考，其中想起我的隐修士朋友们，特别是那位曾教导我、留我、又失去我的神父的模样。我正沉思时，看到一大群蚂蚁在一条小径上蜂拥而过。它们负载的东西明显比自己身体还大。有的用钳子拖着草籽；有的从洞穴里挖土并堆起来以防水。有一队，考虑到冬天临近，并且为了防止储存的谷物因地面潮湿而发芽，正在剪断它们运进来的谷粒的尖端；另一队则庄重地哀悼着搬运死者。更奇怪的是，在这些大军中，出来的并没有阻碍进去的；相反，如果看到有谁被负载压倒，它们就会把肩膀凑上去帮助它。简而言之，那天给了我愉快的娱乐。所以，想起撒罗满如何让我们向蚂蚁的机敏学习，用这样的榜样激励我们迟钝的才能，我开始厌倦被掳，怀念隐修士的小室，渴望模仿那些蚂蚁和它们的行为，在那里劳苦是为了团体，既然没有任何东西属于任何人，所有东西都属于所有人。\par

8. 当我回到卧室，我的妻子迎上来。我的表情显露出内心的忧伤。她问我为何如此沮丧。我把原因告诉她，并劝她逃走。她没有拒绝。我请求她对这件事保持沉默。她答应了。我们不断地低声交谈；在希望与恐惧之间悬而不决。我的羊群中有两只非常好的公山羊：我杀了它们，把它们的皮做成皮囊，用它们的肉准备路上的食物。然后，在傍晚时分，当我们的主人以为我们已经休息时，我们开始了旅程，带着皮囊和一部分肉。当我们到达大约十英里外的河边时，我们吹胀了皮囊，跨坐在上面，把自己托付给水流，用脚慢慢地推动，以便顺流而下到达对岸一个远低于我们出发点的地点，这样追捕者就可能失去踪迹。但与此同时，肉变得湿软，部分丢失了，我们只能靠它维持不到三天的食物。我们尽可能多地喝水，以准备承受预期的口渴，然后匆匆离开，不停地回头张望，夜间比白天走得更快，既因为游荡的撒拉逊人的埋伏，也因为太阳的酷热。即使现在讲述所发生的事，我也感到恐惧；虽然我的心完全平静，但我的身体从头到脚都在颤抖。\par

9. 三天后，我们看到远方有两个骑骆驼的人迅速接近。我立刻预感不祥，开始认为主人打算处死我们，我们的太阳似乎又变暗了。在我们恐惧之中，正当我们意识到沙地上的脚印暴露了我们时，我们在右手边发现了一个向地下延伸很远的洞穴。于是我们进了洞穴：但我们害怕毒兽，如毒蛇、蛇怪、蝎子和其他这类生物，它们常常为了躲避烈日而到这样的阴凉之处。因此我们只勉强进去，躲进左边的一个坑里，不敢再迈一步，以免逃离死亡却陷入死亡。我们心里这样想：如果主在我们苦难中帮助我们，我们就找到了安全；如果他因我们的罪而弃绝我们，我们就找到了坟墓。你认为我们当时是什么感受？我们有多么恐惧，当洞穴前面近处，我们的主人和同伴仆人，被我们的脚印引到我们的藏身之处时？预期的死亡比实际施加的死亡更可怕！我的舌头再次因痛苦和恐惧而结巴；仿佛我听到了主人的声音，几乎不敢嘟囔一句话。他派他的仆人把我们拖出洞穴，他自己则牵着骆驼，手持剑等待我们出来。与此同时，仆人进入了大约三四腕尺，我们在藏身处看到了他的背，但他看不见我们，因为眼睛的特性是：从阳光下进入阴影中的人什么也看不见。他的声音在洞穴中回响：\Quote{出来，你们这些恶棍；出来受死；你们为什么停留？为什么拖延？出来吧，你们的主人在叫你们，耐心地等着你们。}他还在说话时，看哪！透过黑暗，我们看到一头母狮抓住了那个人，勒死了他，把他浑身是血地拖进更深处。好耶稣！那时我们有多么恐惧，又是多么强烈的喜乐！我们看见，虽然我们的主人不知道，我们的敌人灭亡了。他见仆人久久不归，以为逃跑的两个人正在反抗（二对一）。他怒不可遏，仍然手持剑，来到洞穴，像疯子一样大喊，斥责仆人的迟缓，但在他到达我们的藏身处之前，就被野兽抓住了。谁会相信，在我们眼前，一头野兽会为我们而战？\par

一个恐惧的源头被移除，但我们自己也有可能遭遇类似的死亡，尽管狮子的愤怒并不像人的愤怒那样难以忍受。我们因恐惧而心惊胆战：不敢移动一步，等待着结果，在如此巨大的危险中，我们唯一的防御墙就是对我们贞洁的意识；直到清晨，母狮害怕某种陷阱，并意识到自己已被看见，便用嘴叼起幼崽将其带走，留下我们占据我们的避难所。我们的信心并没有一下子恢复。我们没有冲出去，而是等了很长时间；因为每次想到要出去，我们就想象遇到她的可怕情景。\par

10. 最后我们摆脱了恐惧；那天过去后，我们在傍晚时分出发，看到骆驼（因其速度极快而被称为单峰驼）正在安静地反刍。我们骑上它们，借助新获得的谷物的力量，经过十天的沙漠旅行，到达了罗马军营。被呈递给 tribune 后，我们讲述了全部经过，然后被送往在美索不达米亚指挥的萨比亚努斯那里，我们在那里卖掉了骆驼。我亲爱的老院长如今已安眠于主内；我于是来到此地，回到隐修生活，而我把我的同伴托付给修女们照管；因为我爱她如姊妹，却没有把她当作自己的姊妹来对待。\par

玛尔古斯是老人，我是青年，他向我讲述这些事。如今向你们讲述这些事的我已年老，我把它们当作贞洁的历史，为贞洁的人而陈述。童贞女们，我劝勉你们，守护你们的贞洁。将这个故事讲给后来的人，好让他们明白，在刀剑和旷野的野兽中间，德行永远不会被俘，而奉献于基督服务的人可以死，但不能被征服。\par

\SourceRecord{Translated by W.H. Fremantle, G. Lewis and W.G. Martley.；Nicene and Post-Nicene Fathers, Second Series；Vol. 6.；Edited by Philip Schaff and Henry Wace.；Buffalo, NY: Christian Literature Publishing Co.,；1893.；<http://www.newadvent.org/fathers/3006.htm>.}
